% =========================================================================
% CHAPTER 1 — INTRODUCTION
% =========================================================================

\chapter{Introduction}
\label{ch:introduction}

\noindent
{\rmfamily\itshape
\begin{tabular}[t]{@{}l@{\hspace{0.3cm}}p{12.5cm}@{}}
\textls[50]{\textbf{Addressed:}} &
\textls[50]{Background \& Context · Personal Mission Alignment · Thesis Statement · Innovation Relevance \& SDG Alignment}
\end{tabular}
}

% -------------------------------------------------------------------------
\section{The Problem: Cognitive Inaccessibility of Complex Climate Systems}
\label{sec:problem-of-knowing}

For much of the past three decades, climate communication has largely been guided by the \emph{information-deficit hypothesis}: the idea that public inaction on climate change results primarily from a lack of scientific knowledge, and that providing more accurate information would naturally lead to greater understanding and action. However, empirical evidence has consistently challenged this assumption. \citet{kellstedt2008} found, using a representative survey of the American public, that individuals with higher levels of climate knowledge often reported \emph{lower} levels of personal concern and a weaker sense of responsibility for addressing climate change, even after accounting for political affiliation and demographic factors. Similarly, \citet{kahan2015} showed that scientific literacy does not necessarily reduce disagreement. Instead, it can amplify existing cultural divisions, with the most scientifically literate individuals becoming even more polarised when interpreting the same climate evidence. These findings suggest that information alone is rarely sufficient. How people understand climate science depends not only on the information they receive, but also on the cognitive processes, prior beliefs, and cultural contexts through which that information is interpreted. The challenge, therefore, is not simply to communicate climate science more accurately, but to communicate it in ways that better reflect how people perceive, interpret, and make sense of complex scientific information.

This challenge is further complicated by a more fundamental cognitive limitation that predates the climate debate itself. Through two decades of experimental studies, \citet{sterman2008} has shown that even highly educated individuals --- including graduate students at MIT, engineering professionals, and participants with formal training in mathematics and dynamic systems --- often struggle to reason correctly about simple stock-and-flow structures. \citet{cronin2009} demonstrated this limitation through a particularly revealing experiment: participants were presented with the inflow and outflow rates of a system and asked to draw the trajectory of the accumulated stock, a task that requires mentally integrating the difference between two curves. Despite the apparent simplicity of the task, many participants produced incorrect trajectories, confusing rates with levels, interpreting peaks in flows as peaks in accumulated quantities, or reasoning about a dynamic system as if it were static. Importantly, these errors were not limited to participants with low technical ability; they were also observed among individuals with strong mathematical backgrounds who possessed the necessary formal knowledge to solve the problem.

This finding has direct implications for climate communication. Presenting variables such as CO\textsubscript{2} emissions and atmospheric CO\textsubscript{2} concentration as static curves does not automatically convey the underlying relationship between them: the latter represents the accumulation of the former over time. As a result, audiences may systematically underestimate the persistence and inertia of climate processes \citep{sterman2008, cronin2009}. Cognitive-load theory provides a complementary explanation for why these difficulties emerge. Human working memory has a limited capacity for processing and manipulating information, meaning that representations requiring people to mentally reconstruct dynamic processes that are not explicitly shown place a substantial cognitive burden on the viewer. As system complexity increases, so does the likelihood of misunderstanding \citep{sweller1988, swellervanmerrienboer2019}.

Together, these two bodies of literature --- research on how climate information is interpreted and research on how humans reason about dynamic systems --- point towards a common conclusion: communicating complex climate systems is not only a challenge of providing accurate information or improving communication strategies. It is fundamentally a problem of representation. The challenge is to develop forms that allow the dynamic, multi-scale, and non-linear behaviour of climate systems to become perceptually accessible, rather than requiring audiences to reconstruct these processes mentally from static descriptions.


% -------------------------------------------------------------------------
\section{The Gap: Why Conventional Approaches Fall Short}
\label{sec:gap-approaches}

The representational challenge described in the previous section has been studied from different perspectives. Over the past two decades, two largely independent research areas have investigated related aspects of this problem and have produced important insights. The main limitation is not that either field has been unsuccessful, but that their combination --- which is the central idea of this thesis --- has not yet been fully explored.

The first relevant research area is climate-data visualisation. \citet{harold2016} analysed current approaches in climate visualisation through the perspective of cognitive and perceptual science, showing that many common formats used to communicate climate change --- such as temperature trends, projected-change maps, ensemble plots, and uncertainty diagrams --- do not always correspond to how non-expert audiences process visual information. This issue has been supported by further empirical evidence. \citet{mcmahon2015} found that even trained users can have difficulties interpreting uncertainty in IPCC visualisations, often misunderstanding uncertainty ranges and ensemble spread. Similarly, \citet{kause2020} showed that visualisations perceived as clearer, more engaging, or more trustworthy are not necessarily those that lead to better understanding of the underlying scientific concepts. This difference between perceived clarity and actual comprehension highlights the risk of creating visualisations that are attractive but cognitively limited \citep{sheppard2015}. Overall, the literature shows that traditional visual formats often struggle to communicate complex climate systems effectively to non-specialist audiences. However, a systematic alternative that better represents dynamic and complex processes is still missing.

The second relevant research area concerns the use of machine learning in climate and Earth-system science. Reviews by \citet{reichstein2019} and \citet{sonnewald2021} describe the rapid growth of machine-learning applications in areas such as climate forecasting, anomaly detection, physical-process parameterisation, and statistical downscaling \citep{ham2019, weyn2020, rasp2018, vandal2019}. In the context of the AMOC, machine-learning methods have also been applied to investigate variability patterns, early-warning signals, and potential indicators of critical transitions \citep{boers2021, ditlevsen2023, benyami2024}. These approaches have produced valuable scientific results, but their main objective remains predictive: machine learning is primarily used to improve quantitative analysis and modelling for expert users, evaluated through metrics such as prediction accuracy, calibration, and reconstruction performance. Less attention has been given to another possibility: using machine learning to uncover hidden structures in scientific data and translate them into forms that support understanding by non-specialist audiences. This role of machine learning remains largely unexplored in the context of climate tipping systems.

The research gap addressed by this thesis emerges from the intersection of these two areas. Climate visualisation research has identified the limitations of existing representations, but has not fully incorporated the analytical capabilities of modern computational methods. At the same time, machine learning has developed powerful techniques for extracting complex patterns from climate data, but these techniques have rarely been considered as tools for improving scientific understanding and communication.

This thesis addresses this representational challenge by proposing a framework in which unsupervised machine learning is used not primarily as a predictive tool, but as a means of extracting and translating complex system structures into experiential representations. The aim is to achieve enhanced non-sepcialist understanding of complex climate systems by combining machine-learning-derived structures with principles from cognitive science, information visualisation, and embodied cognition. 

The following chapters show that the theoretical foundations and technical components required for such an approach already exist across several disciplines, but have rarely been integrated into a unified framework for communicating complex climate tipping systems. Related approaches have been explored in adjacent fields \citep{iten2020, sacha2018}, but their application to complex climate tipping systems has not yet been systematically investigated. By bringing these perspectives together and validating the approach through the case study of a highly complex, dynamic, and largely invisible Earth-system process (namely the AMOC --- Atlantic Meridional Overturning Circulation), this thesis proposes a methodological contribution at the intersection of machine learning, scientific visualization, cognitive science, and climate communication.

% -------------------------------------------------------------------------
\section{Thesis Statement: An Integrative Framework}
\label{sec:proposal}

This thesis proposes an integrative framework to address the research gap identified in the previous section and to improve the understanding of complex climate systems and tipping risk among non-specialist audiences. The framework combines three complementary stages, each drawing on a different research tradition while providing the foundation for the next.

The first stage is \emph{analytical}. Machine learning is applied to observational and reanalysis data, not to improve prediction, but to recover the latent structure of a complex climate system. The goal is to obtain a compact representation that captures its essential dynamics to be directly use as a driver for more effective forms of representation. The analytical stage is described in detail in Chapter~\ref{ch:data_pipeline}, which focuses on the AMOC as a case study.

The second stage is \emph{design-based}. These latent representations are translated into experiential visualisations using principles from cognitive science, embodied cognition, and information visualisation. Each design decision is explicitly grounded in the relevant cognitiveliterature (ch.~\ref{ch:literature}). A more detailed discussion of the design stage is provided in Chapter~\ref{ch:translate}, which presents the resulting visualisations and explains how they were developed to support human understanding.

The third stage is \emph{empirical}. The resulting visualisations are evaluated against conventional static representations through a mixed-methods user study, investigating how different representational forms influence the understanding of non-specialist audiences. Chapter ~\ref{ch:evaluate} presents the design of the user study, its results, and their implications for the proposed framework.

Together, these three stages form a single methodological pipeline: machine learning extracts latent structures, design transforms them into experiential representations, and empirical evaluation assesses whether these representations improve understanding, compared with conventional static visualisations. 

The framework is organised around one central research question and three subordinate research questions. The central research question asks:

\begin{quote}
\emph{To what extent do machine-learning-driven experiential visualisations improve the cognitive understanding of complex climate tipping systems among non-specialist audiences, compared with conventional static visualisations?}
\end{quote}

This is fundamentally a methodological question. Rather than seeking a new physical discovery or proposing a new machine-learning algorithm, it investigates whether integrating machine learning, cognitive science, and experiential visualisation can improve the communication of complex climate systems.

The three sub-questions mirror the three stages of the framework. The first asks which latent structures can be recovered from observational data through machine learning. The second investigates how these structures can be translated into experiential visualisations grounded in established principles of perception, cognition, and information visualisation. The third evaluates how non-specialist audiences understand complex climate systems when interacting with conventional static visualisations, machine-learning-driven interactive visualisations, and machine-learning-driven experiential visualisations.

Taken together, these questions provide the coordinated evidence required to evaluate the proposed framework.

The main contribution of this thesis is methodological and machine learning serves as a \emph{substrate}, not as a deliverable. Rather than introducing a new climate-science result or a novel machine-learning technique, it proposes and evaluates a framework that connects the latent structure of complex climate systems with forms of representation designed to support human understanding. Its originality lies in the integration of analytical, design, and empirical research within a single coherent workflow. Although demonstrated using the Atlantic Meridional Overturning Circulation (AMOC), the framework is intended to be transferable to other invisible and complex climate systems. The final deliverable is \emph{intelligibility}.

% -------------------------------------------------------------------------
\section{The Case Study: AMOC as a Test of the Framework}
\label{sec:case-study-amoc}

The framework proposed in the previous section requires a case study through which it can be demonstrated. The choice of this case study is therefore not arbitrary: a system used to test whether complex climate dynamics can be made understandable to non-specialist audiences should contain the very characteristics that make communication difficult.\footnote{The AMOC is used a case study. The thesis does not aim to provide new scientific insights into the AMOC itself by making predictions, nor does it contribute to oceanography research, but rather it uses it as a test case for the proposed framework. Therefore, it focuses only on the aspects of the AMOC relevant to the representational challenge addressed by this thesis.}

The Atlantic Meridional Overturning Circulation (AMOC) represents such a case. It is an invisible system in a fundamental sense: it has no direct surface expression that can be observed without instruments, and no everyday experience provides intuitive access to its behaviour. It operates across an entire ocean basin, meaning that its dynamics cannot be understood from a single location. Its variability occurs over timescales of decades to centuries, longer than both individual human experience and the available direct observational record. Moreover, it is a non-linear system and a potential tipping element, where a substantial weakening could lead to a transition towards a different circulation regime with possible hysteresis effects \citep{lenton2008, armstrongmckay2022}. Finally, its behaviour has important climatic implications, as a significant slowdown could alter North Atlantic heat transport and influence the climate of Europe \citep{buckley2016}.\footnote{The physical foundations of the AMOC, including its governing mechanisms, the Stommel two-box model of bistability, and the mathematical formulation of meridional transport, are discussed in Chapter~\ref{ch:background}.}

Physically, the AMOC is the meridional component of large-scale North Atlantic circulation that transports warm surface waters northward and returns colder, denser waters southward at depth. This overturning circulation is driven by density differences related to temperature and salinity variations (thermohaline processes). Its strength is conventionally measured in Sverdrup (\SI{1}{\sverdrup} = \SI{e6}{\cubic\metre\per\second}). Since April 2004, it has been continuously monitored by the RAPID-MOCHA-WBTS observing array at \ang{26.5}N \citep{frajkawilliams2021, moat2026}. Although this record covers only about two decades, while the main modes of AMOC variability occur over multidecadal timescales, it provides the observational basis for the analytical work presented in Chapter~\ref{ch:data_pipeline}. This record is complemented by the Copernicus Marine Service GREP reanalysis ensemble, which extends the analysis beyond the observational period while introducing ensemble spread as an additional representation of uncertainty.

The AMOC is the system studied in this thesis, but it is not the system that most people associate with climate change. When non-specialist audiences are asked to imagine climate change, they are more likely to think about visible phenomena such as melting ice, extreme heat, or rising sea levels \citep{oneill2014}. This difference highlights an important distinction between climate systems that are directly observable and those that require instruments, models, and computational methods to become accessible.

The AMOC belongs to the latter category. However, these two categories are not independent. The Greenland Ice Sheet, one of the most visible indicators of climate change, is also connected to AMOC dynamics. Freshwater released from Greenland enters the North Atlantic, reduces surface water density, weakens deep-water formation, and can contribute to a slowdown of the circulation \citep{boning2016, rahmstorf2015, caesar2018}.\footnote{Although this mechanism is physically well established, the magnitude of Greenland's contribution to recent AMOC changes remains uncertain. Direct RAPID observations do not yet isolate a clear freshwater signal from natural variability, and the relative contribution of Greenland forcing, internal variability, and other external drivers remains an active research question \citep{caesar2018}.}

This connection between the visible and the invisible is central to the argument of this thesis. What people can directly observe --- for example, glacier retreat in Greenland --- is physically linked to a process that remains hidden: the changing circulation of the ocean that regulates heat transport across the North Atlantic. This relationship is used both in the analytical framework of Chapter~\ref{ch:data_pipeline} and in the experiential design developed in Chapter~\ref{ch:translate}.

The need for clearer representations is further reinforced by the fact that even within the scientific community, the recent behaviour of the AMOC remains actively debated. \citet{caesar2018} reconstructed a fingerprint of AMOC weakening from sea-surface temperature patterns and suggested that the circulation had reached an unusually weak state. \citet{boers2021} applied early-warning-signal methods to reconstructed records and reported evidence consistent with a system approaching a critical transition. \citet{ditlevsen2023} further extended this analysis by proposing a probabilistic tipping window centred around the middle of the twenty-first century, a result that received considerable attention but also significant methodological criticism \citep{benyami2024}.\footnote{\citet{benyami2024} questioned the robustness of the statistical assumptions underlying the analysis of \citet{ditlevsen2023}, particularly regarding the use of reconstructed sea-surface temperature patterns as indicators of AMOC stability. At the time of writing, this debate remains unresolved. However, the existence of such disagreement illustrates the difficulty of interpreting and communicating complex climate-system behaviour.}

At the same time, the direct observational record from RAPID has not shown a clear long-term slowdown consistent with some reconstructed estimates \citep{moat2026, rahmstorf2015}. The result is a system with major climatic implications, complex evidence, and dynamics that cannot be directly observed. These characteristics make the AMOC an appropriate test case for investigating whether new forms of representation can improve understanding of complex climate systems.

% -------------------------------------------------------------------------
\section{On the Interdisciplinary Nature of the Work}
\label{sec:positioning}

The work presented here sits at an intersection that reflects the author's cross-disciplinary background and holistic approach to problem solving, bringing together science communication, physical and cognitive sciences, and computational methods. It combines five complementary research fields around a shared objective: improving how non-specialist audiences understand complex climate tipping systems.

Each discipline contributes a different part of the solution. Climate and Earth-system science provides the physical system under investigation, the observational data, and the scientific knowledge that grounds the study. Machine learning provides the analytical methods used to identify latent structures and patterns within high-dimensional climate data. Cognitive science explains why complex dynamic systems are often difficult to understand and identifies the perceptual conditions that can support learning. Information visualisation provides principles for transforming scientific information into meaningful visual experiences, while science communication frames the broader goal of making complex scientific knowledge accessible without compromising its scientific accuracy.

The contribution of this thesis does not lie in any one of these disciplines individually, but in their integration. Scientific evidence, the computational methods used to analyse it, and the way it is ultimately experienced by people are often treated as separate problems. This thesis argues that they should instead be considered parts of the same process. Improving non-expert of climate tipping systems requires not only rigorous scientific analysis, but also representations that reflect how people perceive, interpret, and interact with complex dynamic phenomena.

Within this framework, machine learning is used in a role that differs from its more common predictive applications. Rather than generating forecasts or automating decisions, it is used to uncover latent structures in climate data that can serve as the foundation for new forms of representation. These learned structures become the basis for designing interactive and experiential visualisations that make otherwise invisible system dynamics more perceptually accessible. While related ideas have begun to emerge in neighbouring fields \citep{iten2020}, their application to climate communication remains largely unexplored. This thesis investigates that possibility using the Atlantic Meridional Overturning Circulation (AMOC) as a case study, with the broader goal of developing a transferable framework for communicating complex climate tipping systems.

% -------------------------------------------------------------------------
\section{Contribution, Innovation Relevance and SDG Alignment}
\label{sec:contribution}

As already mentioned above, the contribution of this thesis is methodological. It does not introduce a new finding about the ocean, nor a new machine-learning technique. Instead, it proposes a defensible framework for connecting the analytical structure of a complex climate system with the perceptual conditions through which non-specialists encounter it, together with an empirical evaluation of whether this connection supports understanding. The field addressed by this work can be situated within the emerging area of \emph{complexity communication}, which investigates how complex, dynamic and often invisible systems can be made cognitively accessible. The contribution does not lie in the novelty of any single component, but in their integration within a coherent research programme and in their evaluation against a conventional baseline, allowing the effectiveness of the approach to be assessed through evidence rather than assumption.

The novelty of the proposed approach becomes clearer when positioned against existing research, which can be organised into four main families. The first concerns conventional static visualisations of climate data, whose limitations for non-specialist audiences have been widely documented \citep{harold2016, kause2020}; these representations constitute the baseline that this thesis seeks to extend rather than replace. The second concerns interactive climate-data visualisation, including institutional tools such as the University of Maine's Climate Reanalyzer \citep{climatereanalyzer} and the Copernicus Interactive Climate Atlas \citep{c3satlas}. Although these platforms provide valuable access to climate information, they primarily rely on conventional charts and maps, and machine learning, where used, generally supports prediction rather than the construction of new representational forms. The third family concerns immersive and experiential approaches to scientific communication, which have demonstrated the potential of interactive, embodied and sensory formats, including through empirical evaluation of audience responses \citep{zelada2024unnatural}. However, these approaches have rarely been built upon machine-learning-derived representations of scientific systems and have not specifically addressed climate tipping dynamics. The fourth family is data physicalisation, which explores how data can become a material and embodied experience \citep{lupi2016}, but has not been extensively applied to complex climate systems. No existing approach combines the four elements brought together in this thesis: the recovery of latent structures from a climate tipping system through machine learning, their translation into experiential representations, their grounding in cognitive-science principles, and their comparative empirical evaluation against conventional visual formats. The contribution therefore lies in this specific integration and in demonstrating a methodology through which its effectiveness can be investigated.

The relevance of this work operates across three complementary levels, which this thesis aims to substantiate without extending beyond the evidence it can provide. 

The first level is \emph{academic relevance.} Research on the communication of complex climate systems remains relatively limited in methodological terms. Existing studies have often focused on visual preferences, asking participants which representations they perceive as clearer, more engaging or more trustworthy, while direct assessment of conceptual understanding remains less common \citep{harold2016}. Experiential representations, including interactive, embodied or sensory approaches, are still largely absent from systematic evaluation frameworks. This thesis contributes not another visualisation alone, but a comparative methodology in which the same underlying scientific information is translated into three representational conditions of increasing experiential richness and evaluated through qualitative investigation of the understanding they support. To the author's knowledge, this approach has not yet been systematically applied to the AMOC or to a climate system with comparable structural complexity. 

The second level is \emph{institutional relevance}, where the alignment with the United Nations Sustainable Development Goals is substantive rather than nominal. The primary connection is with \textbf{SDG 13 (Climate Action)}, and specifically target 13.3, which calls for improved education, awareness-raising and capacity building on climate change \citep{un2015}. These objectives directly correspond to the central aim of this thesis: improving how complex climate risks are communicated and understood. A secondary connection is with \textbf{SDG 4 (Quality Education)}, as the work addresses the challenge of making scientifically complex concepts accessible to non-specialist audiences. A further, indirect connection is with \textbf{SDG 14 (Life Below Water)}, given that the AMOC is part of the ocean processes that influence marine and coastal systems. The alignment therefore does not rely on thematic overlap alone, but on the contribution of the thesis to the educational and awareness-building goals identified by these frameworks.

The third level is \emph{practical relevance}, concerning the potential transferability of the proposed methodology beyond the specific case study. The AMOC represents one example of a broader class of climate tipping elements, including the Greenland and West Antarctic ice sheets, the Amazon rainforest, boreal permafrost, the Indian summer monsoon and coral reef systems, which share several characteristics that make them difficult to communicate: limited direct perceptibility, long timescales, non-linear behaviour, potential threshold dynamics and significant societal consequences \citep{armstrongmckay2022, lenton2008}. 

The framework developed in this thesis — machine learning as an analytical substrate, experiential representation as a communicative layer, and qualitative evaluation as a feedback mechanism — is designed with transferability in mind. The AMOC therefore represents a challenging test case rather than the final application of the approach. If supported by the evidence collected in this thesis, the methodology could provide a foundation for communicating other complex climate systems whose scientific importance currently exceeds their public legibility. These three levels reinforce one another: the academic gap motivates the methodological contribution, the institutional relevance supports the choice of case study, and the practical dimension justifies exploring a framework that may extend beyond a single system. Before a society can act on a complex system, it must first be able to form a meaningful understanding of it; this thesis contributes to the ongoing effort of making that understanding possible.

% -------------------------------------------------------------------------
\section{Roadmap}
\label{sec:roadmap}

The thesis is organised in three parts, each addressing one stage of the detect–translate–evaluate arc that structures the argument.

\textbf{Part~I --- Foundations} establishes the problem this thesis addresses and the case study through which it is investigated. Chapter~\ref{ch:rqs} formalises the research questions, objectives and hypotheses that structure the work. Chapter~\ref{ch:literature} reviews the relevant literature across five domains --- psychological barriers to climate risk perception, the cognitive science of dynamic systems, empirical evaluation of climate visualisation, embodied and immersive representation, and machine-learning applications to climate science --- and identifies the intersection at which this thesis positions itself. Chapter~\ref{ch:background} develops the physical and mathematical apparatus needed to treat the AMOC as a measurable and analysable object.

\textbf{Part~II --- Translation} documents the constructive work of the thesis. Chapter~\ref{ch:methodology} specifies the three-stage methodology --- analytical, design-based and empirical --- through which the argument is made testable. Chapter~\ref{ch:data_pipeline} presents the analytical pipeline: the dimensionality-reduction, anomaly-detection and early-warning analyses through which machine learning recovers the latent structure of the observational record, and reports the results that motivate the representational choices that follow. Chapter~\ref{ch:translate} designs and documents the three representational conditions --- conventional static, interactive dynamic and machine-learning-driven experiential --- through which that latent structure is translated for non-specialist audiences.

\textbf{Part~III --- Evaluation and Reflection} closes the argumentative loop. Chapter~\ref{ch:evaluate} specifies the qualitative user study through which the three conditions are compared. Chapter~\ref{ch:results} reports its findings. Chapter~\ref{ch:discussion} interprets those findings in conversation with the literature reviewed in Part~I and with the analytical results of Part~II, and Chapter~\ref{ch:conclusion} concludes with reflections on limitations, transferability to other invisible climate systems and future work.
